@DATABASE PGS.TEX
@(C) 1995 Soft-Logik Publishing Corporation
@INDEX PageStream3:Help/PGS.INA/MAIN



@remark *** TEXT SUBJECTS HELP FILE ***

@node ENTR "Entering Text"

@{FG Highlight}Entering Text@{FG Text}

You can type text directly in PageStream or you can @{"import" link IMPO} it from a word
processor. PageStream can import the text from a word processor file and
maintain its text attributes.

PageStream also offers a wide range of word processing features to allow
you to enter and edit text for your document without having to use a word
processor. You may prefer to type directly in PageStream for small blocks
of text. If you have a faster computer, you may wish to type all your
text directly in PageStream.

@{FG Highlight}Typing Text@{FG Text}

After you have created a text frame (text column, text object or text
shape), you can enter text in it by typing directly into PageStream, or
by @{"importing" link IMPO} a text file created by another program. You must position
the text @{"insertion point" link INSP} before you can enter text. The insertion point
is a vertical line that indicates where the new text will appear. As you
type, text appears at the insertion point as the insertion point moves to
the right.

If your text reaches the right margin as you continue to type, it will
wrap down to the next line automatically; you do not need to press Return
to end a line as you would on a typewriter. Use the Return key only when
you want to end a paragraph.

See also @{"Typing Special Characters" link TYPS}.

@{FG Highlight}Ending a Paragraph@{FG Text}

Paragraphs are ended by pressing the Return key. Ending a paragraph
forces the text to wrap to the next line without continuing to the right
edge. Paragraph spacing is normally used to set the space between
paragraphs, but you can also press the Return key a second time to insert
a blank line between paragraphs.

If you end a paragraph by holding down a Shift key while pressing Return,
PageStream will ignore the paragraph spacing value and will start the
next paragraph on the next line.

You can break a paragraph into two paragraphs by moving the cursor to the
desired position and pressing the Return key.

@{FG Highlight}Correcting Mistakes@{FG Text}

When you make a mistake, you can correct it by backspacing over it with
the Backspace key. You can erase characters anywhere in your text by
positioning the insertion point in it and pressing the Backspace or Del
keys. The Backspace key deletes the character to the left of the
insertion point, and the Del key deletes the character to the right of
the insertion point.

If the insertion point is at the beginning of a paragraph, pressing
Backspace will erase the end of paragraph of the previous paragraph,
joining the two paragraphs together. Similarly, if the insertion point is
at the end of a paragraph, pressing Del will erase the end of paragraph,
joining the two together.

@{FG Highlight}Overset Text@{FG Text}

If there is too much text to fit into an article's frame(s), the text
will be overset. Overset frames can be distinguished by the overset
indicator at their bottom right corner. (Refer to page 164 of the manual
for an example of an overset indicator.)

You can search for overset text with the @{"Find Overset Text" link PageStream3:Help/PGS.sub/OVER} macro.
@toc PageStream3:Help/PGS.HEL/MAIN
@endnode

@node INSP "Insertion Point"

@{FG Highlight}Insertion Point@{FG Text}

You must select the @{"Text" link PageStream3:Help/pgs.too/TEXT} tool before you can enter new text or edit
existing text in a text column or object. When the Text tool is selected,
the mouse pointer will change to the text cursor shape.

You must position the text insertion point before you can enter text. The
insertion point is a vertical line that indicates where the new text will
appear. As you type, text appears at the insertion point as the insertion
point moves to the right.

@{FG Highlight}Placing the Insertion Point@{FG Text}

If you choose the Text tool when a text frame is selected, the insertion
point will be placed automatically in that frame. If one was not
selected, or if you wish to enter text in a different text frame, you
must place the insertion point manually. To place the insertion point in
a text frame, move the mouse cursor over it at the desired point and
click the left mouse button. If the frame does not yet contain text, the
insertion point will be placed at its top left corner. If there is text
at the point where you click, the insertion point will be positioned
there.

@{FG Highlight}Moving the Insertion Point@{FG Text}

You can move the insertion point anywhere in your text to edit it in
three ways:

- position the insertion point with the mouse cursor
- use the arrow keys
- use a keyboard shortcut

To move the insertion point with the mouse cursor, move the mouse so that
the mouse cursor is between the characters where the insertion point
should be positioned. Click the mouse the place the insertion point at
the mouse position.

To move one character left or right, press the left or right arrow key.
To move up or down one line, press the up or down arrow key. You can hold
down arrow keys to move the insertion point repeatedly. You can also use
the arrow keys in conjunction with the Ctrl, Alt and Shift keys to move
the cursor greater distances.
@toc PageStream3:Help/PGS.HEL/MAIN
@endnode

@node TYPS "Typing Special Characters"

@{FG Highlight}Typing Special Characters@{FG Text}

PageStream allows you to enter many special characters that are not shown
on the keyboard. The KeyShow program in your Tools directory shows you
how to type some of these special characters. Run KeyShow and select the
Alt key or the Alt and Shift keys to see which characters are available
using standard Amiga keystrokes.

PageStream allows you to enter more characters than provided by the Amiga
operating system, so it offers another way to enter special characters.
Each character has a character identification number, and a few of the
common ones also have mnemonic sequences. To enter a character by its
character identification number, simply press Ctrl-D and then type the
unicode character number. To enter a character by its mnemonic sequence,
press Ctrl-C and then type the mnemonic.

Some common typographic characters and their identification numbers are:
@{code}
@{FG Highlight}Special Character   Identification Number (Ctrl-D)  Mnemonic (Ctrl-C)@{FG Text}
left double quote   08220                           "l
right double quote  08221                           "r
bullet              08226                           bu
em dash             08211                           m-
en dash             08212                           n-
trademark           08482                           tm
copyright mark      00169                           co
registered mark     00174                           rm
paragraph marker    00182                           p!
section marker      00167                           se
@{body}
@{body}
The @{"Insert Character" link PageStream3:Help/pgs.typ/INCH} command can be used to enter a character without
knowing its identification number or mnemonic.
@endnode

@node IMPO "Import/Export Filters"

@{FG Highlight}Import/Export Filters@{FG Text}

PageStream's @{"Insert Text" link PageStream3:Help/PGS.PRO/INTX} command allows you to import text created with a
word processor or text editor. You can import text only from programs for
which PageStream has an import/export filter, and from programs that
allow you to save in ASCII format. Import/export filters translate text
from word processor formats into a format that PageStream can understand
and are also used for exporting text to word processing formats.

The text import/export filters provided with this release of PageStream3
are:

@{FG Highlight}Filters@{FG Text}
ASCII.tfilter
FinalWriter.tfilter
Wordworth.tfilter
Excellence.tfilter
ProWrite.tfilter
WordPerfect.tfilter
Word.tfilter
IFFCTXT.tfilter
IFFDTXT.tfilter
IFFFTXT.tfilter

Soft-Logik makes every attempt to ensure complete compatibility with the
word processors listed. However, because they are often updated and
changed, Soft-Logik cannot guarantee complete compatibility.
@toc PageStream3:Help/PGS.HEL/MAIN
@endnode

